%; whizzy paragraph -pdf xpdf -latex ./whizzypdfptex.sh
%; whizzy-paragraph "^\\\\begin{frame}"
% latex beamer presentation.
% platex, latex-beamer でコンパイルすることを想定。 

%     Tokyo Debian Meeting resources
%     Copyright (C) 2009 Junichi Uekawa
%     Copyright (C) 2010 Nobuhiro Iwamatsu

%     This program is free software; you can redistribute it and/or modify
%     it under the terms of the GNU General Public License as published by
%     the Free Software Foundation; either version 2 of the License, or
%     (at your option) any later version.

%     This program is distributed in the hope that it will be useful,
%     but WITHOUT ANY WARRANTY; without even the implied warranty of
%     MERCHANTABILITY or FITNESS FOR A PARTICULAR PURPOSE.  See the
%     GNU General Public License for more details.

%     You should have received a copy of the GNU General Public License
%     along with this program; if not, write to the Free Software
%     Foundation, Inc., 51 Franklin St, Fifth Floor, Boston, MA  02110-1301 USA

\documentclass[cjk,dvipdfmx,12pt]{beamer}
\usetheme{Tokyo}
\usepackage{monthlypresentation}
%  preview (shell-command (concat "evince " (replace-regexp-in-string "tex$" "pdf"(buffer-file-name)) "&"))
%  presentation (shell-command (concat "xpdf -fullscreen " (replace-regexp-in-string "tex$" "pdf"(buffer-file-name)) "&"))
%  presentation (shell-command (concat "evince " (replace-regexp-in-string "tex$" "pdf"(buffer-file-name)) "&"))

%http://www.naney.org/diki/dk/hyperref.html
%日本語EUC系環境の時
\AtBeginDvi{\special{pdf:tounicode EUC-UCS2}}
%シフトJIS系環境の時
%\AtBeginDvi{\special{pdf:tounicode 90ms-RKSJ-UCS2}}

\title{東京エリア Debian 勉強会}
\subtitle{資料}
\author{岩松 信洋 iwamatsu@debian.or.jp\\IRC nick: iwamatsu}
\date{2010年01月23日}
\logo{\includegraphics[width=8cm]{image200607/openlogo-light.eps}}

\begin{document}

\frame{\titlepage{}}

\emtext{設営準備にご協力ください}


\begin{frame}
 \frametitle{2010年01月}
\begin{minipage}[t]{0.45\hsize}
  \begin{itemize}
  \item 注意事項
	\begin{itemize}
	 \item 飲食OK
	\end{itemize}
  \item 最近あったDebian関連のイベント報告
	\begin{itemize}
	 \item 前回の勉強会
	\end{itemize}
 \end{itemize}
\end{minipage} 
\begin{minipage}[t]{0.45\hsize}
 \begin{itemize}
  \item 09:00 - BSP 開会の挨拶 / 荒木 靖宏
  \item 09:10 - BSP についての説明 / 岩松 信洋
  \item 09:20 - 12:00 BSP 午前の部 / 参加者全員
  \item 12:00 - 13:00 昼食タイム
  \item 13:00 - 13:30 BSP 中間発表 / 参加者全員
  \item 13:30 - 19:00 BSP 午後の部 / 参加者全員
  \item 19:00 - 19:30 BSP 結果発表 / 参加者全員
  \item 19:30 - 後片付け / 参加者全員
 \end{itemize}
\end{minipage}
\end{frame}

\section{}
\begin{frame}
 \frametitle{Agenda}
\begin{minipage}[t]{0.45\hsize}
  \begin{itemize}
  \item 注意事項
	\begin{itemize}
	 \item 飲食可能
	\end{itemize}
  \item 最近あったDebian関連のイベント報告
	\begin{itemize}
	 \item 前回の勉強会
	 \item Adam C Powell IV 来日
	 \item Dirk Eddelbuettel 来日
	\end{itemize}
 \end{itemize}
\end{minipage} 
\begin{minipage}[t]{0.45\hsize}
 \begin{itemize}
  \item 一年の勉強会を振り返る
  \item 2009年、2010年、2011年
  \item qemubuilder
  \item lxc
  \item BSP, BSP, BSP
 \end{itemize}
\end{minipage}
\end{frame}

\begin{frame}{Hack Cafe}

毎週水曜日、週に一回東京のどっかのカフェでハック。\\
\url{http://twitter.com/debian_hackcafe}\\
関西でも始めたようです。
\end{frame}

\emtext{事前課題}

\begin{frame}{事前課題}
\begin{enumerate}
 \item  今回の BSP への意気込みを熱く語ってください。
\end{enumerate}
\end{frame}

{\footnotesize
\begin{prework}{$B$d$^$@(B}
\preworksection{$B:#2s$N(B BSP $B$X$N0U5$9~$_$rG.$/!)8l$C$F$_$k(B}

Primer$B$H(BRC-bug$B%j%9%HFI$_$J$,$i05E]$5$l$D$D$b!"(B

\begin{itemize}
\item $B$b$&!V$$$D=P$k$N$+K\Ev$K=P$k$N$+!W$H$O8@$o$;$J$$!*(B
\item $B$d$C$Q$j!V%^%@!<!)!W$h$j<+J,$G2?$+$7$F=P$kJ}$,$$$$$h$M!*(B
\item $B$d$k$J$i%H%3%H%s=8CfE*$K$d$C$F%a%s%F%J%9%-%k8~>e$b$7$h$&!*(B
\end{itemize}

$B$H$$$&$3$H$G!"(B10$B8D0LD>$9$3$H$rL\I8$K;22C$7$^$9!#(B

$B<B$O!V=i;22C$,$$$-$J$j(BBSP$B$C$FBg>fIW$+!)!d<+J,!W$J$N$G(B
$B1[$($i$l$J$$JI$r$$$^C[$$$F$7$^$C$?46$,L5$-$K$7$bHs$:$G$9$,!"(B
$BLnK>C#@.$N$?$a$;$C$;$H;E9~$s$G$f$-$^$9$G$9!&!&!&(B

\end{prework}

\begin{prework}{$B2,It(B $B5f(B}
\preworksection{$B:#2s$N(B BSP $B$X$N0U5$9~$_$rG.$/8l$C$F$/$@$5$$(B}
$B$^$:$O<j;}$A$N%Q%C%1!<%8$G$"$k(Buim$B$N(Blintian$B%P%0$r(B0$B$K$7$^$9!#(B
$B$=$N$"$H!"(Buim$B$N%S%k%@$+$i$N%(%i!<Js9p$N2r(B $Bat(B $BO$r$d$j$?$$$H;W$$$^$9!(B?

$B$=$l$G$bM>NO$,$"$C$?$i(Bmlterm$B$N(Blintian$B%P%0$d$=$NB>$N%j%j!<%9%/%j%F%#%+%k(B
$B%P%0$rD>$=$&$H;W$$$^$9!#(B
\end{prework}

\begin{prework}{ $BF|HfLn(B $B7<(B }

groff $B$N%P%0$r<h$k$N$KD)@o$7$F$_$h$&$H;W$$$^$9!#(B

$B$"$H4X?t7?$^$o$j$G2?$+$"$l$P6(NO$7$^$9!#(B

\end{prework}



\begin{prework}{ $BF#BtM}Ao(B(risou) }

$B2?$r$7$?$iNI$$$N$+$bA4A3$o$+$C$F$^$;$s$,!"$3$N5!2q$rMxMQ$7$F!"(BDebian$B$K>/(B
 $B$7$G$b9W8%$G$-$l$P!"$H;W$C$F$$$^$9!#$h$m$7$/$*4j$$$7$^$9!#(B

\end{prework}



\begin{prework}{ kmuto }

$B<+J,$N%Q%C%1!<%8$N%P%0=$@5!"(BRC$B$K8B$i$:;H$C$F$F:$$k%Q%C%1!<%8$N(Bdelay
 upload$B!"$[$+$N?M$N:n6H$N%9%]%s%5!<!"$H$$$C$?$"$?$j$r!#(B
$B;~4|E*$K$A$g$C$HFI$a$J$$$N$G%-%c%s%;%k$9$k$+$b$7$l$^$;$s!#$"$H!";22C$7$F(B
 $B$bESCfB`@J$K$J$k$H;W$$$^$9!#(B


\end{prework}



\begin{prework}{ henrich }

$B<j;}$A$N:n6H$r>/$7$G$b8:$i$7$F$*$3$&$H;W$$$^$9!#(B

\end{prework}



\begin{prework}{ iwamatsu }
RC$B$H(BQA$B$rCf?4$K%P%0$rDY$9M=Dj$G$9!#(B


\end{prework}



\begin{prework}{ $B>e@n=c0l(B }

$B%,%s%,%s%P%0$r$D$V$9M=Dj$G$9!#(B
$B$"$H!"#D#e#b#i#a#nJY6/2qM=Ls%7%9%F%`$K$D$$$F>pJs6&M-$7$?$$$G$9!#(B

\end{prework}

}

\emtext{BSP 開会の挨拶}
\begin{frame}{BSP 開会の挨拶}
Debian Developer である荒木さんから、開会の挨拶と本年度の目標等について
 ありがたいお言葉を頂戴します。
\end{frame}

\emtext{BSP についての説明} 

\begin{frame}{BSPとは}
バグを潰しまくるパーティです。バグ増やさないようにしましょう。
\end{frame}

\begin{frame}{RC バグの対応}
\begin{enumerate}
 \item RCバグを直したい人は \#debian-bugs にログインしてください。
 \item \#debian-bugs で 対象のバグを直している人がいないか確認します。
 \item \url{http://wiki.debian.org/BSP2010/CoordinationTokyoMoenchengladbach}
      にRCバグを登録してください。
 \item バグを直します。
 \item その辺にいる Debian Developerに相談して、NMUしてもらいます。 
\end{enumerate}
\end{frame}

\emtext{BSP 午前の部}
\emtext{BSP 中間発表}
\emtext{BSP 午後の部}
\emtext{BSP 成果発表}
\emtext{後片付け}

\begin{frame}{次回の勉強会}

\begin{itemize}
 \item 2010年2月: 木更津開催
 \item 2010年2月: OSC Tokyo Spring
\end{itemize}
 
\end{frame}

\end{document}

;;; Local Variables: ***
;;; outline-regexp: "\\([ 	]*\\\\\\(documentstyle\\|documentclass\\|emtext\\|section\\|begin{frame}\\)\\*?[ 	]*[[{]\\|[]+\\)" ***
;;; End: ***
